\begin{resumo}
  A síntese de imagens realísticas é um problema central da Computação Gráfica,
  em que se busca reproduzir algoritmicamente o transporte de luz de forma
  fisicamente plausível. O \emph{Path Tracing} é uma das abordagens que realiza
  menos simplificações desse transporte, permitindo que efeitos como reflexões
  e iluminação indireta surjam naturalmente do algoritmo. Por estimar a equação
  da renderização \citep{kajiya86} por meio de um integrador de Monte Carlo, no
  entanto, o método é extremamente custoso computacionalmente, e sua qualidade
  depende diretamente do número de amostras utilizadas. Neste trabalho,
  implementamos um \emph{Path Tracer} e exploramos a Amostragem de Importância
  Múltipla \citep{veach95}, uma técnica que combina diferentes estratégias de
  amostragem de forma a reduzir a variância do estimador. A implementação
  combina amostragem da distribuição de refletância e amostragem direta de
  fontes de luz, e a derivação dos pesos é apresentada de forma explícita para
  o caminho completo. A efetividade da técnica é avaliada por meio da
  comparação de imagens renderizadas com igual número de amostras, com e sem a
  Amostragem de Importância Múltipla. Os resultados demonstram que a técnica
  reduz significativamente a variância do estimador: imagens de variância
  renderizadas com igual número de amostras mostram uma redução expressiva, com
  a versão sem Amostragem de Importância Múltipla apresentando variância
  visivelmente elevada em contraste com a versão que a utiliza.

\end{resumo}
